\section{Introduction}
Many garbage-collected languages, such as Ruby, Python, JavaScript and Java, are implemented in low-level languages (e.g., C and C++) that do not have garbage collection mechanisms.  In order for garbage collectors (GCs) to correctly recognize live objects, object manipulations in the low-level implementations need to accompany with ``GC operations'' that inform GCs of object references stored and passed around in registers and memory.

Inserting those operations is critical yet error-prone.  If an implementation omits a GC operation, GC can reclaim objects that are still in use or cannot reclaim dead objects.  At the same time, as GC operations are ``extra'' to the core implementation logic, omission of a GC operation merely affects GC behaviors, which can be unnoticed for a long time.  For example, in the reference implementation of Ruby (CRuby)\footnote{\url{https://github.com/ruby/ruby}}, there are new bugs that omit extra operations for GC even after the implementation has been maintained 30 years\cite{newly_found_missing}.

Statically ensuring that an implementation correctly has GC operations is desired. This is because dynamic measurements such as testing are not suitable due to non-deterministic and hard-to-reproducible nature of GC errors.  In fact, there has been an attempt to statically check insertion of GC operations in a compact implementation of a garbage-collected language for embedded systems\cite{UgawaFujimoto2019PRO}.

We present GuardLint, a static analyzer for CRuby that checks GC operations in its implementation. The analysis is challenging because of the large codebase of CRuby (more than 300 thousand lines of code), the complicated condition when a GC operation is needed due to the CRuby's conservative GC, ``subordinate memory areas,'' and low-level compiler optimizations (explained in Section~\ref{sec:Background}), and the fact that excessive GC operations (i.e., performing GC operations where they are not needed) can cause performance degradation (Section~\ref{sec:Benchmark}).

In the rest of the paper, we first review CRuby's memory management system, including the GC operation, called \texttt{RB\_GC\_GUARD} (Section~\ref{sec:Background}). We then analyze the challenges of statically detecting guard omissions in a codebase of this scale (Section~\ref{sec:Challenges}). Section~\ref{sec:Approach} presents our analyzer's design and implementation. Section~\ref{sec:Experiment} reports the results of the analyzer and additional manual inspection. Section~\ref{sec:Discussion} discusses limitations and future work. Section~\ref{sec:Related} reviews related work. Section~\ref{sec:Conclusion} concludes the paper.
